Use the theta characteristic $\kappa$ to centre the theta divisor on
$J(C)=\operatorname{Pic}^0(C)$:
\[
 \Theta_\kappa=\{M\in J(C):h^0(C,\kappa\otimes M)>0\}.
\]
Let $L=\mathcal O_{J(C)}(\Theta_\kappa)|_P$.  The type of the restricted
polarization can be read off directly.  Under the principal polarizations of
the two Jacobians, $\pi^*$ and $\Nm_\pi$ are dual homomorphisms; hence the
orthogonal complement of $P$ in $J(C)$ is $\pi^*J(E)$
\cite[Section~12.3]{BirkenhakeLange2004}.  Therefore
\[
 \ker(\varphi_L)=P\cap\pi^*J(E).
\]
Since $\pi^*$ is injective and $\Nm_\pi\circ\pi^*=[3]$ on $J(E)$,
\[
 P\cap\pi^*J(E)=\pi^*J(E)[3].
\]
This group has order $9$.  If the polarization type is
$(d_1,d_2,d_3)$ with $d_1\mid d_2\mid d_3$, then
$|\ker(\varphi_L)|=(d_1d_2d_3)^2$; hence $d_1d_2d_3=3$ and necessarily
\begin{equation}
 \operatorname{type}(L)=(1,1,3),
 \qquad L^3=3!\,d_1d_2d_3=18.
\label{eq:prym-polarization}
\end{equation}
Scheme-theoretically,
\[
 X=\Theta_\kappa\cap P
  =\{M\in P:h^0(C,\kappa\otimes M)>0\}.
\]
The theta divisor is preserved as a closed subscheme by the Serre-duality
involution $L'\mapsto K_C\otimes(L')^{-1}$ on $\operatorname{Pic}^3(C)$.
Since $\kappa^2=K_C$, this involution is $M\mapsto M^{-1}$ in centred
coordinates.  Relative Serre duality preserves the determinantal ideal of
theta, so after restriction to $P$ one obtains
\begin{equation}
 X=-X
\label{eq:symmetry}
\end{equation}
as closed subschemes.  The centre will be seen to be unique in
\eqref{eq:stabilizer}.

\subsection{The norm fibre and the resolution}

Let
\[
 n:C^{(3)}\longrightarrow\operatorname{Pic}^3(E),
 \qquad D\longmapsto\mathcal O_E(\pi_*D),
\]
and put
\[
 N=n^{-1}(\eta^3),
 \qquad
 a:N\longrightarrow X,
 \qquad
 D\longmapsto\mathcal O_C(D)\otimes\kappa^{-1}.
\]
For the centred Abel map $D\mapsto\mathcal O_C(D)\otimes\kappa^{-1}$ one has
scheme-theoretically
\[
 N=C^{(3)}\times_{J(C)}P.
\]
Thus $a$ is its base change to $P$.  The map $a$ is surjective: if $M\in X$,
a nonzero section of $\kappa\otimes M$ has an effective divisor $D$ of
degree three, and
\[
 \Nm_\pi\mathcal O_C(D)
 =\Nm_\pi(\kappa\otimes M)=\eta^3.
\]
The finite morphism $C^{(3)}\to E^{(3)}$ induced by $\pi$ restricts to a
finite surjection
\[
 N\longrightarrow|\eta^3|\simeq\PP^2.
\]
Hence $\dim N=2$, so the restriction of the centred theta section to $P$ is
nonzero.  Therefore $X\in|L|$ is an effective Cartier divisor, and it is
ample.

Let $\Gamma=|\kappa|\simeq\PP^1$ and put
\[
 D_0=r_1+r_2+r_3\in\Gamma.
\]

\begin{lemma}[The norm fibre]\label[lemma]{lem:norm-fibre}
The surface $N$ has the unique singular point $D_0$, and the germ $(N,D_0)$
is a simple elliptic singularity of type $\wt E_6$.  If
$\beta:\widehat N\to N$ is the blowup of $D_0$, then $\widehat N$ is smooth,
its exceptional curve $E_0$ is elliptic with $E_0^2=-3$, and the strict
transform $\widetilde\Gamma$ is a $(-1)$-curve meeting $E_0$ transversely in
one point.  Contracting $\widetilde\Gamma$ gives a smooth surface $S$ and a
proper morphism
\[
 \rho:S\longrightarrow X.
\]
The image $E_{\exc}$ of $E_0$ satisfies $E_{\exc}^2=-2$ and maps to the
origin of $X$.
\end{lemma}

\begin{proof}
Because the three branch points are distinct, a neighbourhood of $D_0$ in
$C^{(3)}$ is the product of neighbourhoods of the $r_i$.  Choose additive
analytic coordinates $t_i$ on $E$ at $b_i$ so that the Abel map
$E^{(3)}\to\operatorname{Pic}^3(E)$ is, after translation, given by
$t_1+t_2+t_3$.  After rescaling local parameters $z_i$ at $r_i$, the cover
has $t_i=z_i^3$.  Hence the norm fibre has the analytic equation
\[
 z_1^3+z_2^3+z_3^3=0.
\]
Thus $(N,D_0)$ is the affine cone over a smooth plane cubic, hence a simple
elliptic singularity of type $\wt E_6$.

It remains to locate the other critical points of the norm map.  Separate
the possible multiplicity patterns on the symmetric cube.  A divisor
containing an unramified point
has a tangent direction on which $d\pi$ is nonzero, so it cannot be critical.
Near a ramification point the map is $z\mapsto z^3$.  If two points collide,
then in the elementary symmetric coordinates $e_1=z_1+z_2$ and
$e_2=z_1z_2$ one has
\[
 z_1^3+z_2^3=e_1^3-3e_1e_2,
\]
which has no linear term at the double ramification cycle.  If three points
collide, however,
\[
 z_1^3+z_2^3+z_3^3=e_1^3-3e_1e_2+3e_3,
\]
and the $3e_3$ term makes the differential nonzero.  Thus, apart from the
cycle $D_0=r_1+r_2+r_3$, the only possible critical cycles are
$2r_i+r_j$ with $i\ne j$.  Membership in the norm fibre would then force
\[
 2b_i+b_j\sim3O.
\]
If $k$ is the third index, \eqref{eq:branch-sum} gives
$b_i+b_j+b_k\sim3O$, hence $b_i=b_k$, contradicting the fact that the branch
points are distinct.  Therefore $D_0$ is the only singular point of $N$.

Blowing up the vertex of the cubic cone gives a smooth surface $\widehat N$
with exceptional elliptic curve $E_0$ and $E_0^2=-3$.  The local behaviour
of the pencil $\Gamma=|\kappa|$ at the vertex is as follows.  Projection
formula gives
\[
 \pi_*\kappa\simeq\eta\oplus\mathcal O_E\oplus\eta^{-1}.
\]
The two sections of $\kappa$ may therefore be chosen as follows: the section
coming from the $\mathcal O_E$-summand is represented by $w$ and has zero
divisor $D_0$, while a generator of $H^0(E,\eta)$ pulls back to a section
that is nonzero at each $r_i$, since none of the branch points $b_i$ is $O$.  Let $t$ be the affine parameter of the pencil
centred at the member $D_0$.  In a local trivialization of $\kappa$ at
$r_i$, the moving section has equation
\[
 w+t u_i=0,\qquad u_i(r_i)\ne0.
\]
Since $w$ is a local parameter at the totally ramified point $r_i$, after
rescaling $z_i$ this gives
\[
 z_i(t)=-u_i(r_i)t+O(t^2).
\]
Thus $\Gamma\to N$ has a nonzero tangent vector at $D_0$ and hence
multiplicity one there.  Because $\Gamma\subset N$, its tangent direction
lies on the exceptional cubic
\[
 \lambda_1^3+\lambda_2^3+\lambda_3^3=0,
 \qquad \lambda_i=-u_i(r_i),
\]
and the displayed expansion shows that the strict transform of $\Gamma$
meets $E_0$ in one reduced point, transversely.

Let $a_C:C^{(3)}\to J(C)$ be the Abel--Jacobi map and let $\Theta_C$ be the
principal theta divisor.  If $X_{p_0}\subset C^{(d)}$ denotes the divisor of
cycles containing a fixed point $p_0$, Macdonald's canonical-class formula is
\cite{Macdonald1962}
\[
 K_{C^{(d)}}\equiv a_d^*\Theta_C+(g-d-1)X_{p_0}.
\]
For $g(C)=4$ and $d=3=g-1$ the second term vanishes, hence
\[
 K_{C^{(3)}}\equiv a_C^*\Theta_C.
\]
Since $N$ is a fibre of $n:C^{(3)}\to\operatorname{Pic}^3(E)$, its normal
line bundle is trivial.  Adjunction therefore gives
\[
 K_N\equiv a^*L.
\]
The map $a$ contracts $\Gamma$, so $K_N\cdot\Gamma=0$.  The cubic cone has
discrepancy $-1$.  Here $N$ is a Cartier divisor in the smooth variety
$C^{(3)}$, hence is Gorenstein.  The numerical equivalence above is used
only for intersection numbers, while the local cubic-cone equation gives
the divisor identity
\[
 K_{\widehat N}=\beta^*K_N-E_0.
\]
Since $E_0\cdot\widetilde\Gamma=1$, one has
$K_{\widehat N}\cdot\widetilde\Gamma=-1$.  Adjunction on
$\widetilde\Gamma\simeq\PP^1$ then gives
$\widetilde\Gamma^2=-1$.  Contracting it produces a smooth surface $S$ and
\[
 E_{\exc}^2=E_0^2+(E_0\cdot\widetilde\Gamma)^2=-2.
\]
The composite $\widehat N\to N\to X$ is constant on
$\widetilde\Gamma$, hence factors through $\rho:S\to X$.  The curve
$E_{\exc}$ maps to the origin.
\end{proof}
